\chapter{Sprint 4.1 --- Fix BA spec gaps + Báo cáo + Pop-up chọn cửa hàng}

\section{Bối cảnh}

Sau khi Sprint 4 push lên GitHub, review lại \emph{2 - FE Portal} sheet → phát hiện 6 chỗ lệch:

\begin{itemize}
    \item 4 bug số đếm dashboard không đúng spec.
    \item 2 trang Phase 1 BA bắt buộc bị xếp nhầm vào Phase 2: \emph{Báo cáo đặt hàng} + \emph{Pop-up chọn cửa hàng}.
    \item Tab \emph{Lịch sử thông báo} chưa tách khỏi \emph{Thông báo từ HQ}.
\end{itemize}

Sprint 4.1 fix toàn bộ trong cùng commit, giữ tinh thần \emph{follow đúng BA Phase 1}.

\section{4 bug fix dashboard}

\begin{longtable}{p{4.5cm} p{4.5cm} p{5.5cm}}
\toprule
\textbf{Block dashboard} & \textbf{BA yêu cầu} & \textbf{Đã fix} \\
\midrule
\endhead
Return requests count & state in (\texttt{sent},\texttt{approved}) — KHÔNG tính \texttt{draft} & Loại \texttt{draft} ra khỏi list. \\
Most purchased products & \texttt{date\_order >= today - 90 days} & Thêm filter 90 ngày trong \texttt{\_top\_products}. \\
Latest notifications + Latest returns & \texttt{limit 3} & Đổi từ 5 → 3. \\
Latest returns list & \texttt{state != cancel} (skeleton: \texttt{!= rejected}) & Thêm \texttt{('state', 'not in', ['rejected'])}. \\
\bottomrule
\end{longtable}

Verify: \texttt{anh.owner} có 1 RR \texttt{sent} + 1 RR \texttt{done} → counter dashboard hiển thị \emph{1} (chỉ sent).

\section{Module \texttt{wujia\_portal\_report} (BA Phase 1)}

URL \texttt{/portal/reports/orders}, role check: chỉ \emph{Owner} hoặc \emph{Manager} truy cập, \emph{Staff} bị 303 redirect.

\subsection{4 block render trong 5 query}

\begin{itemize}
    \item \textbf{Block 1 --- 4 KPI cards:} Tổng đơn / Doanh thu / Đơn hoàn thành / Đơn hủy. Từ \emph{1 \texttt{\_read\_group}} group by \texttt{state} → tách 4 KPI từ kết quả.
    \item \textbf{Block 2 --- Bar chart 12 tháng:} đơn theo tháng + doanh thu (column + line). Dùng \texttt{\_read\_group(groupby=['date\_order:month'])}.
    \item \textbf{Block 3 --- Top 10 sản phẩm:} \texttt{sale.order.line.\_read\_group} group by \texttt{product\_id}, sum qty + total. \emph{1 query}.
    \item \textbf{Block 4 --- Donut state distribution:} reuse data Block 1 (không query thêm). Render donut + table tỷ lệ \%.
\end{itemize}

\subsection{Performance design}

\begin{lstlisting}[style=python]
# 5 query / page render. Pre-format số liệu ở controller.
state_groups = SO._read_group(domain=base, groupby=['state'],
                              aggregates=['__count', 'amount_total:sum'])
# → 4 KPI từ 1 result, donut chart từ cùng result.

month_groups = SO._read_group(domain=..., groupby=['date_order:month'], ...)
top_groups = SOL._read_group(domain=..., groupby=['product_id'], limit=10, ...)
# Chart payload dump JSON 1 lần qua data-chart-payload attribute,
# ApexCharts client đọc + render, KHÔNG AJAX thêm.
\end{lstlisting}

Default date range = năm hiện tại — giới hạn dataset, dùng index \texttt{date\_order} + \texttt{franchise\_id} đã có.

\section{Pop-up chọn cửa hàng (BA Phase 1)}

User multi-franchise như \texttt{dung.multi} (3 cửa hàng) khi vào \texttt{/portal} sẽ thấy \emph{modal chặn} (Bootstrap \texttt{static} backdrop) yêu cầu chọn 1 cửa hàng. Sau khi chọn, cookie \texttt{wujia\_active\_franchise\_id} được set 30 ngày.

\subsection{Cookie pattern (ADR-021)}

Lưu \texttt{active\_franchise\_id} vào \emph{cookie HTTP} thay vì \emph{DB column trên \texttt{res.users}}:

\begin{itemize}
    \item Cookie: max-age 30 ngày, samesite=Lax, httponly=False (JS đọc tên cửa hàng).
    \item Validate ở controller: cookie \emph{phải in} \texttt{user.\_get\_accessible\_franchise\_ids()} (ormcached). Cookie không hợp lệ → fallback all.
    \item Single-franchise user → auto-pick first (không cần modal).
\end{itemize}

\textbf{Lý do không dùng DB column:}
\begin{itemize}
    \item DB write mỗi lần switch → tốn I/O.
    \item Phải invalidate ormcache \texttt{\_get\_accessible\_franchise\_ids} khi đổi.
    \item Cookie expire tự nhiên → user idle 30+ ngày → reset mặc định, đỡ stale.
    \item Cookie cùng scope với session, logout xóa kèm.
\end{itemize}

\subsection{UX flow}

\begin{lstlisting}[style=shell]
1. dung.multi login → / portal → must_pick_franchise=True → modal HIỆN.
2. Chọn cửa hàng → POST /portal/franchise/switch → cookie set → 303 redirect.
3. F5 dashboard → modal KHÔNG hiện. Top navbar badge "Đang tại: <tên>" hiện.
4. Click badge → modal hiện lại (data-action="open-store-picker").
5. Logout → cookie xóa kèm session.
\end{lstlisting}

\subsection{Helper module-level cho các module portal\_*}

\texttt{wujia\_portal\_base/controllers/portal.py} export 3 function:

\begin{itemize}
    \item \texttt{get\_active\_franchise\_id()} --- đọc cookie + validate, cache trong request lifetime.
    \item \texttt{get\_active\_franchise\_ids\_filter()} --- tuple cho \texttt{('franchise\_id', 'in', ...)} domain.
    \item \texttt{get\_max\_role\_in\_franchises()} --- role cao nhất user có (Owner > Manager > Staff).
\end{itemize}

9 callsite trong 7 module portal\_* refactor từ \texttt{user.\_get\_accessible\_franchise\_ids()} → \texttt{get\_active\_franchise\_ids\_filter()}. Khi chưa chọn cửa hàng → fallback all (giữ behavior cũ).

\section{Tab "Mới" / "Lịch sử thông báo"}

Trong \texttt{/portal/notification}, thêm 2 tab Bootstrap với badge số đếm:

\begin{itemize}
    \item Tab \emph{Mới}: \texttt{date >= today - 30 days}.
    \item Tab \emph{Lịch sử}: \texttt{date < today - 30 days}.
\end{itemize}

Cùng 1 controller, query string \texttt{?tab=recent|history}. Domain riêng, KHÔNG thêm route.

\section{ADR mới}

\subsection{ADR-021: Cookie HTTP cho \texttt{active\_franchise\_id}}
Đã giải thích ở section pop-up. Trade-off chính: cookie rẻ hơn DB column khi user chuyển cửa hàng thường xuyên, đặc biệt với 1500 user concurrent.

\subsection{ADR-022: Module-level helper cho portal\_* modules}
3 function module-level (\texttt{get\_*}) trong \texttt{wujia\_portal\_base/controllers/portal.py} — các module phụ \texttt{from odoo.addons.wujia\_portal\_base.controllers.portal import ...}. Pattern này đơn giản hơn extend \texttt{res.users} vì cookie cần \texttt{request} (model-level method với \texttt{ormcache} không thấy được request).

\section{Verify Sprint 4.1 (đã pass)}

\begin{lstlisting}[style=shell]
14 routes test với 3 user role:
  - anh.owner (owner, 1 cửa hàng): tất cả 200, badge "Đang tại: Cầu Giấy"
                                    auto hiện vì single franchise.
  - dung.multi (manager, 3 cửa hàng): tất cả 200, modal chặn dashboard
                                       cho đến khi chọn. Cookie set sau switch.
  - cuong.staff (staff): /portal/reports/orders → 303 redirect (role check).

Bug R1 verify: anh.owner có 1 RR sent + 1 RR done.
  - Dashboard counter "Đổi trả đang mở" hiển thị 1 (chỉ sent, KHÔNG tính done).

Tab notification: cnt_recent=10, cnt_history=0 (noti demo đều mới).
\end{lstlisting}

\section{Files thay đổi Sprint 4.1}

\begin{lstlisting}[style=shell]
SỬA:
  custom/wujia_portal_base/controllers/portal.py     # 4 bug fix + 3 helper + 2 route
  custom/wujia_portal_base/__manifest__.py           # bump v19.0.4.0.0 + thêm files
  custom/wujia_portal_layout/views/layouts.xml       # giữ nguyên (inherit ở base)
  custom/wujia_portal_notification/controllers/      # tab Mới/Lịch sử
  custom/wujia_portal_notification/views/            # 2 tab + badge số đếm
  custom/wujia_portal_notification/static/src/css/   # tab styling
  custom/wujia_portal_*/controllers/portal.py        # 7 file: refactor 9 callsite
  docs/wujia-tea-doc.tex                              # chapter Sprint 4.1

TẠO MỚI:
  custom/wujia_portal_base/views/store_picker_modal.xml
  custom/wujia_portal_base/views/store_picker_navbar.xml
  custom/wujia_portal_base/static/src/js/store_picker.js
  custom/wujia_portal_base/static/src/css/store_picker.css
  custom/wujia_portal_report/                         # module mới (KPI + chart + role)
\end{lstlisting}

% ===========================================================================
